\section{Certified clusters: subspaces, not individual states}
\label{sec:subspace}

The certificate of Section~\ref{sec:cert} answers the question ``which
eigenvalue is this state?''. At a degenerate level that question has no
answer: the problem does not define a single eigenvector, and a network that
solves it exactly returns an arbitrary element of a two-dimensional space
that depends on its seed. The honest object to certify is therefore the
\emph{invariant subspace}, and the honest statement is about the cluster of
eigenvalues it spans. This section gives that statement; it is what turns the
refusal of Section~\ref{sec:degeneracy} into a positive result.

Guaranteed two-sided bounds for eigenvalue clusters, multiplicity included,
already exist for conforming discretizations \citep{cances2020}. The setting
here differs in one respect that changes what has to be computed: a network
output is not a Galerkin projection, so no discretization parameter controls
its error and no saturation hypothesis relates it to an enriched space. Both
the cluster interval and the multiplicity must therefore be obtained from the
count, and the vector-side statement is a principal-angle bound that assumes
nothing about how the trial subspace was produced.

\subsection{The trial subspace and its residual}
\label{sec:sub}

Let $\psi_1,\dots,\psi_m\in\R^{|I|}$ be $m$ approximations---$m$ network
outputs, $m$ snapshots of one training run, or $m$ vectors from any other
source---and let $X=[x_1,\dots,x_m]$ be an orthonormal basis of their span,
$\langle x_i,x_j\rangle_h=\delta_{ij}$, obtained by symmetric
orthogonalization $X=\Psi G^{-1/2}$, $G_{ij}=\langle\psi_i,\psi_j\rangle_h$.
This choice matters: $G^{-1/2}$ is invariant under rotations of the input
basis, so the certificate depends on $\operatorname{span}\{\psi_j\}$ and not
on how the columns were ordered or mixed. If $G$ is singular the $m$ vectors
do not span an $m$-dimensional space and the certificate refuses before
computing anything else.

The compression of $\Hh$ to the trial subspace and its residual are
\begin{equation}
  M = X^{\mathsf T}\!\Hh X\in\R^{m\times m},\qquad
  R = \Hh X - X M ,
  \qquad \rho_F=\norm{R}_F ,
  \label{eq:rr}
\end{equation}
and the Rayleigh--Ritz pairs $(\theta_i,y_i)$, $My_i=\theta_iy_i$ with
$\theta_1\le\cdots\le\theta_m$, give the Ritz vectors $x_i\leftarrow X y_i$
(we reuse the symbol, since $\{Xy_i\}$ is again an orthonormal basis of the
same subspace) and the residual
\begin{equation}
  r_i=\Hh x_i-\theta_i x_i = R y_i ,\qquad
  \rho_i=\norm{r_i}_h,\qquad \sum_i \rho_i^2=\rho_F^2 .
  \label{eq:ritzres}
\end{equation}
Everything in \eqref{eq:rr}--\eqref{eq:ritzres} is computed from the
approximations alone.

\subsection{Two propositions}
\label{sec:props}

Write $N(\sigma)=\#\{i:\lambda_i<\sigma\}$ for the verified count of
Section~\ref{sec:counting}, and suppose an interval $(a,b)$ has been
\emph{verified} to contain exactly $m$ eigenvalues,
\begin{equation}
  N(b)-N(a)=m ,
  \label{eq:cluster}
\end{equation}
counting multiplicity. Let $V\subset\R^{|I|}$ be the corresponding invariant
subspace, $\dim V=m$, and let $P_V$ be the orthogonal projector onto it in
$\langle\cdot,\cdot\rangle_h$. Define the cluster margin
\begin{equation}
  \delta_i=\min(\theta_i-a,\,b-\theta_i),\qquad \delta=\min_i\delta_i .
  \label{eq:margin}
\end{equation}

\begin{lemma}[Cluster localization]\label{lem:cluster}
  Assume \eqref{eq:cluster}, $\theta_i\in(a,b)$ for all $i$, and
  $\rho_i<\delta_i$. Then
  \[
    \operatorname{dist}\big(\theta_i,\ \spec(\Hh)\cap(a,b)\big)\le\rho_i
    \qquad\text{for every }i .
  \]
\end{lemma}
\begin{proof}
  The pair $(x_i,\theta_i)$ is an approximate eigenpair in the sense of
  Definition~\ref{def:pair} with $\norm{x_i}_h=1$ and residual $r_i$, so
  Lemma~\ref{lem:weinstein} gives
  $\operatorname{dist}(\theta_i,\spec(\Hh))\le\rho_i$. Every eigenvalue of
  $\Hh$ outside $(a,b)$ is at distance at least $\delta_i$ from $\theta_i$,
  and $\rho_i<\delta_i$; hence the eigenvalue realizing the distance lies
  inside $(a,b)$.
\end{proof}

Lemma~\ref{lem:cluster} is the statement the single-state certificate cannot
make: it is valid no matter how close the $\theta_i$ are to each other,
because it never asks the interval to separate them. What it does not give is
an identification: when $\rho_i\ge\delta_i/2$ the intervals
$[\theta_i-\rho_i,\theta_i+\rho_i]$ may cover the same eigenvalue twice, and
the certificate makes no claim about which cluster eigenvalue belongs to
which Ritz value. At a degenerate level no such claim exists.

The angle between the trial subspace and $V$ obeys a Davis--Kahan inequality
in which the role of the spectral gap is played by the margin
\eqref{eq:margin}.

\begin{theorem}[Principal angle]\label{thm:angle}
  Under the hypotheses of Lemma~\ref{lem:cluster},
  \[
    \sin\Theta\big(\operatorname{span}X,\,V\big)\ \le\ \frac{\rho_F}{\delta},
  \]
  where $\Theta$ is the largest principal angle between the two subspaces.
\end{theorem}
\begin{proof}
  Let $\{\varphi_j\}$ be an orthonormal eigenbasis of $\Hh$ with
  $\Hh\varphi_j=\lambda_j\varphi_j$, and expand
  $x_i=\sum_j c_{ij}\varphi_j$. Since $\Hh x_i-\theta_i x_i
  =\sum_j(\lambda_j-\theta_i)c_{ij}\varphi_j$,
  \[
    \rho_i^2=\sum_j(\lambda_j-\theta_i)^2c_{ij}^2
    \ \ge\ \delta_i^2\!\!\sum_{j:\,\lambda_j\notin(a,b)}\!\! c_{ij}^2
    \ =\ \delta_i^2\,\big\lVert P_V^{\perp}x_i\big\rVert_h^2 ,
  \]
  because an eigenvalue outside $(a,b)$ differs from $\theta_i\in(a,b)$ by at
  least $\delta_i$. The Ritz vectors form an orthonormal basis of
  $\operatorname{span}X$, so
  \[
    \big\lVert P_V^{\perp}X\big\rVert_F^2
    =\sum_i\big\lVert P_V^{\perp}x_i\big\rVert_h^2
    \le\sum_i\frac{\rho_i^2}{\delta_i^2}
    \le\frac{\rho_F^2}{\delta^2},
  \]
  and the sine of the largest principal angle is the spectral norm of
  $P_V^{\perp}$ restricted to $\operatorname{span}X$, which is bounded by the
  Frobenius norm.
\end{proof}

Theorem~\ref{thm:angle} is the quantitative statement that a set of network
outputs spans the right subspace. For a single state ($m=1$) it reduces to
the Davis--Kahan bound \eqref{eq:dk} with $\delta$ in place of $\gamma$, which
is the same inequality shifted by $\eta$; the cluster version is not a
sharper theorem, it is the same theorem applied to the object that remains
well defined when the eigenvalue is degenerate.

\subsection{The cluster certificate}
\label{sec:cluster-alg}

The certificate is assembled exactly as in Section~\ref{sec:isolating}, with
$m$ in place of $1$:

\begin{enumerate}
\item \emph{Rank.} Form $G$ and refuse if it is numerically singular.
\item \emph{Rayleigh--Ritz.} Compute $M$, $\theta_i$, $\rho_i$, $\rho_F$
  from \eqref{eq:rr}--\eqref{eq:ritzres}.
\item \emph{Interval.} Propose $a$ and $b$ at the midpoints of the spectral
  gaps on either side of the block of $m$ approximate eigenvalues that
  matches $\{\theta_i\}$ (ARPACK is used to propose, never to decide), then
  \emph{verify} \eqref{eq:cluster}. If the proposal fails, bracket the two
  counting transitions by bisection and verify again; if that fails too, the
  refusal reason distinguishes a multiplicity mismatch from a subspace that
  spans more than one cluster.
\item \emph{Accept} if $\delta$ exceeds the resolution floor \eqref{eq:floor}
  and $\rho_i<\delta_i$ for every $i$; then report the interval $(a,b)$, the
  per-value enclosures $[\theta_i-\rho_i,\theta_i+\rho_i]$, and the angle
  bound of Theorem~\ref{thm:angle}. Otherwise refuse, with the reason.
\end{enumerate}

The cost is two verified counts---one sparse factorization each---one
$m\times m$ symmetric diagonalization, and $m$ applications of $\Hh$; it does
not grow with the number of requested eigenvalues beyond the $m$ matrix--vector
products, and it is negligible next to training.

Three properties are worth stating explicitly. First, $m$ is \emph{measured},
not asserted: the count decides whether the trial subspace has the dimension
the invariant subspace has, and a subspace that missed a state is refused
rather than certified with a smaller dimension. Second, Lemma~\ref{lem:cluster}
brackets each of the $m$ eigenvalues to within $\rho_i$ of a Ritz value, and the
multiplicity itself is verified only down to the counting resolution
\eqref{eq:floor}; both hold independently of how accurate the network is: a
poor subspace yields a poor angle bound but never a wrong multiplicity. Third, the interval may
contain a block of adjacent distinct eigenvalues; degeneracy is the case that
motivates the construction, not a special case of it.

The refusal reasons are meant to be actionable, and that suggests a
procedural use of the certificate. A multiplicity mismatch says that the trial
subspace missed a state of the cluster or covered more than one cluster, and
tells the user how many vectors to supply or how to split them;
$\rho_i\ge\delta_i$ says that the subspace is not an accurate enough
approximation and that more training is required; a span whose second
direction is the difference of nearly identical outputs says that the runs
collapsed and that a further direction must be forced, which is done by
deflating the accepted outputs as well and retraining. The loop
\emph{train, certify, deflate, retrain} is what produced the accepted neural
pair of Section~\ref{sec:deg-results}, and its termination argument is the one
familiar from deflated eigensolvers: each accepted direction is projected out
of the residual that the next run minimises.

\begin{remark}[What a cluster certificate does not say]
\label{rem:cluster-not}
It does not say that $\theta_i$ approximates $\lambda_i$ in the matching
order, and it does not give an $O(\rho^2)$ enclosure: for $m=1$ the
Kato--Temple bound of Theorem~\ref{thm:temple} is sharper and should be
preferred whenever its hypothesis holds. It also does not resolve eigenvalues
that the counting cannot separate, exactly as in \eqref{eq:floor}; what it
does is to make the resolution floor the \emph{only} obstacle, removing the
structural one.
\end{remark}

\begin{remark}[Do independent runs span a degenerate level?]\label{rem:two}
It is tempting to certify a degenerate level by training $m$ networks and
certifying their span, and the answer depends on the data rather than on the
construction. In our experiments three of four independent runs collapsed onto
the same symmetry-broken representative (pairwise overlap $>0.9999$) while the
fourth found the orthogonal partner. The certificate accepted every pair
containing the fourth run, with bounds of $5.6$--$6.3\times10^{-3}$; it refused
the two remaining pairs of near-identical outputs, $(0,1)$ and $(1,2)$, whose
second directions are differences of norm $2.4\times10^{-3}$ and
$6.7\times10^{-3}$---for both, the Ritz residual of that direction ($16.3$ and
$7.5$) exceeds the margin ($6.1$ and $7.2$), so the span contains no invariant
subspace of dimension two; and it accepted the pair $(0,2)$, which omits the
fourth run, only with a bound two orders of magnitude looser ($0.68$), which
is the honest signal that its second direction was mostly noise. The certified
object is the subspace, not the individual vector, and the certificate sorts
sets of outputs by how much of a second direction they actually contain. When the answer is negative the
second direction can be forced by deflating the first output as well, as in
Section~\ref{sec:deg-results}.
\end{remark}
